% §8.6 Bibliographic Notes
%
% Monograph register: the entropy manuscript, the counting argument, entropy and
% capacity, compressed Bloom filters. The compressed-Bloom note now points to the
% corrected rem:compressed-bloom (the earlier version had the mechanism inverted).
% Citations towell-bernoulli-entropy, carter-wegman-1979, cover-thomas,
% shannon-1948, mitzenmacher-2002, mitzenmacher-upfal. Cross-refs thm:entropy-floor,
% prop:bloom-overhead, sec:entropy-floor, sec:entropy-refresher, sec:bsc-capacity,
% rem:compressed-bloom, ch:optimal-bernoulli-sets. DEFINED: sec:bib-notes-ch8.

\section*{Bibliographic Notes}
\addcontentsline{toc}{section}{Bibliographic Notes}
\label{sec:bib-notes-ch8}

\paragraph{The entropy manuscript.}
The floor of \Cref{thm:entropy-floor} and the Bloom overhead of
\Cref{prop:bloom-overhead} follow \textcite{towell-bernoulli-entropy}, the
primary source for \Cref{sec:entropy-floor} and \Cref{sec:bloom-gap}: the
counting argument, the asymptotic $-n\log_2 \fprate$ form, and the $\log_2 e$
comparison are all from it, as is the honest statement that the two-sided case
remains open. This chapter adds the BSC bridge of \Cref{sec:bsc-capacity} that
ties the floor to Part~I's running channel, the entropy refresher for readers
new to $H(X)$, and the explicit forward connection to the chapter-14
construction that closes the gap.

\paragraph{The counting argument.}
The proof of \Cref{thm:entropy-floor} is the classical lower bound of Carter,
Floyd, Gill, Markowsky, and Wegman, in the combinatorial style of the universal
hashing program of \textcite{carter-wegman-1979}: reason about a hash-based data
structure by counting collisions over an explicit family rather than appealing
to a random oracle. For every representation $R$, count the accepting sets it
can describe; for every target set, count the representations that work; the
ratio is the bound.

\paragraph{Entropy and capacity.}
The background of \Cref{sec:entropy-refresher} and \Cref{sec:bsc-capacity} is
standard; \textcite{cover-thomas} covers entropy, mutual information, the BSC,
capacity, and the source-coding theorem in its first three chapters, with the
BSC capacity $C = 1 - H(\fprate)$ and the source-coding bound that underwrites
\Cref{thm:entropy-floor} treated rigorously. \textcite{shannon-1948} remains the
cleanest short presentation of the source-coding idea.

\paragraph{Compressed Bloom filters.}
\textcite{mitzenmacher-2002} introduced the compressed Bloom filter, whose
mechanism \Cref{rem:compressed-bloom} works out: the half-full optimal array is
incompressible, so compression instead exploits a deliberately sparse array,
recovering per-bit redundancy and approaching the floor for the
\emph{transmission} setting it targets, while the in-place query structure the
rest of the chapter analyzes keeps the $\log_2 e$ gap. It is the right entry
point for how much of that gap compression alone can close, and sets up the
contrast with the structural change of \Cref{ch:optimal-bernoulli-sets}.
\textcite{mitzenmacher-upfal} is the companion for the probabilistic background.
